\subsection{Extern verbleibende Prüfungen}
\label{subsec:extern-verbleibende-pruefungen}

Die offenen Punkte in Tabelle~\ref{tab:external-verification-items} sind keine
einheitliche Fehlerklasse. \emph{FAIL} bedeutet, dass eine ausgeführte
Prüfung ihr Kriterium nicht erfüllte; \emph{nicht verifiziert} bezeichnet
fehlende hinreichende Evidenz; \emph{außerhalb des Scopes} kennzeichnet eine
bewusste Produktgrenze. So sind die aktuellen roten CI-Jobs reale
Fehlerbefunde, während ein nie gestarteter Compose-Verbund nicht als
fehlgeschlagen gelten kann.

Der aktuelle Core-Containerjob belegt zwar Compose-Validierung und den Bau
beider Master-Images, aber weder \texttt{docker compose up} noch eine
produktive Bereitstellung. Linux- und macOS-Pakete sind auf
\texttt{a4487ac} unsigniert gebaut; Windows ist wegen des fehlgeschlagenen
Installationsschritts nicht paketiert. Branch Protection konnte über die
öffentliche API ohne Berechtigung nicht gelesen werden. Signing,
Notarisierung, Updater, produktive Zugangsdaten sowie Authentifizierung/RBAC
benötigen schließlich ein konkretes Produkt und geschützte externe
Infrastruktur \parencite{kleiveist2026ci-core-run,kleiveist2026ci-desktop-run,kleiveist2026release}.

\input{workspace/tables/external-verification-items}
